\usepackage{unicode-math}
\setmathfont{texgyrepagella-math.otf}
\usepackage{pgfplots,tocloft}
\usetikzlibrary{shapes.geometric}
\pgfplotsset{compat=1.18}
\setlength{\cftsecnumwidth}{3.3em}
\setlength{\cftsubsecindent}{4.8em}
\setlength{\cftsubsecnumwidth}{4.2em}
\definecolor{Blue}{HTML}{173E60}
\definecolor{Green}{HTML}{26714D}
\definecolor{Red}{HTML}{AA354C}
\definecolor{Pale}{HTML}{EDF4F8}
\definecolor{Gold}{HTML}{F2D799}
\tikzset{flow/.style={-{Stealth[length=2mm]},thick,Blue},
 note/.style={font=\small\sffamily,align=center},
 smallnote/.style={font=\footnotesize\sffamily,align=center}}
\newcommand{\Generated}{}
\newcommand{\PlateDirectory}{}
\newcommand{\ChapterNamespace}{}
% Plates keep their designed type size: wide plates shrink to the measure;
% narrow plates are enlarged by at most 15 percent rather than stretched.
% varwidth measures the plate in paragraph mode, so a figure file may stack a
% picture above a legend line exactly as in its standalone chapter.
\usepackage{varwidth}
\newsavebox{\PlateBox}
% Plates sized from \linewidth still see the real measure inside the wide box.
\newlength{\PlateMeasure}
\newcommand{\PlaceScaledPlate}{%
\ifdim\wd\PlateBox>\linewidth\resizebox{\linewidth}{!}{\usebox{\PlateBox}}%
\else\ifdim1.15\wd\PlateBox<\linewidth\scalebox{1.15}{\usebox{\PlateBox}}%
\else\resizebox{\linewidth}{!}{\usebox{\PlateBox}}\fi\fi}
\newcommand{\Plate}[3]{\par\medskip\noindent\begin{minipage}{\linewidth}\centering
\setlength{\PlateMeasure}{\linewidth}%
\sbox{\PlateBox}{\begin{varwidth}{8\PlateMeasure}\setlength{\linewidth}{\PlateMeasure}%
\centering\input{\PlateDirectory/#1.tex}\end{varwidth}}\PlaceScaledPlate
\captionof{figure}{#2}\begingroup
\edef\CurrentBookFigureLabel{\noexpand\OriginalBookLabel{\ChapterNamespace#3}}\CurrentBookFigureLabel
\endgroup\end{minipage}\par\medskip}
\newcommand{\Listing}[1]{\par\Needspace{5\baselineskip}
\VerbatimInput[breaklines,breakanywhere,fontsize=\small,frame=single,numbers=left,numbersep=4pt]{\Generated/#1.py}\par\medskip}
\let\OriginalBookLabel\label
\let\OriginalBookRef\ref
\let\OriginalBookAutoref\autoref
\makeatletter
% fvextra/lineno can leave a one-column bookkeeping marker in an incremental
% auxiliary file even when the current wrapper does not activate lineno.
% Consume that legacy marker without changing numbered VerbatimInput listings.
\providecommand\@LN@col[1]{}
\let\OriginalBookDisplayLabel\label@in@display
\newcommand{\ScopedChapterLabels}{
\renewcommand{\label}[1]{\edef\CurrentBookLabel{\noexpand\OriginalBookLabel{\ChapterNamespace##1}}\CurrentBookLabel}
\renewcommand{\label@in@display}[1]{\edef\CurrentBookLabel{\noexpand\OriginalBookDisplayLabel{\ChapterNamespace##1}}\CurrentBookLabel}
\renewcommand{\ref}[1]{\OriginalBookRef{\ChapterNamespace##1}}
\renewcommand{\autoref}[1]{\OriginalBookAutoref{\ChapterNamespace##1}}}
\makeatother
\fvset{breaklines=true,breakanywhere=true}
\fancyhead[R]{\small\sffamily Chapter \thechapter}
\setcounter{tocdepth}{1}
\AtBeginDocument{\raggedright\setlength{\parindent}{0pt}}
% ---------------------------------------------------------------------------
% Book-level reading apparatus (prepublication revision, October 2026).
% Chapter sources are unchanged; the continuous builder injects these boxes.
\usepackage{tcolorbox}\tcbuselibrary{skins,breakable}
% tcolorbox loads the verbatim package, which redefines {verbatim}; restore
% the textbook preamble's breaking fvextra environment.
\DefineVerbatimEnvironment{verbatim}{Verbatim}{breaklines,fontsize=\footnotesize}
\definecolor{Ink}{HTML}{1F2A33}
\definecolor{GreenPale}{HTML}{EAF4EE}
\definecolor{GoldPale}{HTML}{FBF4E2}
\definecolor{Slate}{HTML}{5B6B78}
\newtcolorbox{EvolutorGoalsBox}{enhanced,breakable,colback=Pale,colframe=Blue,
 boxrule=0pt,leftrule=2.6pt,arc=0pt,outer arc=0pt,left=9pt,right=9pt,top=6pt,bottom=5pt,
 before skip=10pt,after skip=12pt,
 before upper={{\sffamily\bfseries\small\color{Blue}In this chapter you will\par}\smallskip}}
\newtcolorbox{EvolutorIdeasBox}{enhanced,colback=GreenPale,colframe=Green,
 boxrule=0pt,leftrule=2.6pt,arc=0pt,outer arc=0pt,left=9pt,right=9pt,top=6pt,bottom=5pt,
 before skip=12pt,after skip=12pt,
 before upper={{\sffamily\bfseries\small\color{Green}Key ideas\par}\smallskip}}
\newenvironment{ChapterGoals}{\begin{EvolutorGoalsBox}\small\begin{itemize}
 \setlength{\itemsep}{1.5pt}\setlength{\parskip}{0pt}\setlength{\topsep}{0pt}}%
 {\end{itemize}\end{EvolutorGoalsBox}}
% Key ideas are short; the box is unbreakable and moves whole to the next page.
\newenvironment{KeyIdeas}{\begin{EvolutorIdeasBox}\small\begin{itemize}
 \setlength{\itemsep}{2pt}\setlength{\parskip}{0pt}\setlength{\topsep}{0pt}}%
 {\end{itemize}\end{EvolutorIdeasBox}}
% Generic boxes for front and back matter.
\newtcolorbox{EvolutorNote}[1]{enhanced,breakable,colback=GoldPale,colframe=Gold!80!black,
 boxrule=0pt,leftrule=2.6pt,arc=0pt,outer arc=0pt,left=9pt,right=9pt,top=6pt,bottom=5pt,
 before skip=10pt,after skip=10pt,
 before upper={{\sffamily\bfseries\small\color{Ink}#1\par}\smallskip}}
% Part pages: number, title, guiding question, orientation and chapter list.
\newcommand{\PartBodyFile}{}
\newcommand{\PartRoman}[1]{\uppercase\expandafter{\romannumeral#1}}
\newcommand{\PartAfter}{\vspace{6mm}{\color{Blue}\rule{\linewidth}{0.8pt}}\par\vspace{7mm}%
 \normalfont\normalsize\color{black}\raggedright
 \OriginalBookLabel{book:part\arabic{part}}\input{\PartBodyFile}}
\titleformat{\part}[display]{\raggedright\sffamily\color{Blue}}%
 {\large\bfseries PART \thepart}{6pt}{\Huge\bfseries}[\PartAfter]
\newcommand{\EvolutorPart}[3]{\renewcommand{\PartBodyFile}{#3}\part{#2}}
\titlespacing*{\part}{0pt}{-6mm}{0pt}
\newcommand{\PartQuestion}[1]{{\large\itshape\color{Blue}#1\par}\vspace{3mm}}
\newenvironment{PartChapters}{\par\vspace{5mm}{\sffamily\bfseries\small\color{Blue}CHAPTERS IN THIS PART\par}\vspace{1.5mm}\small}{\par}
\newcommand{\PartChapter}[3]{\par\noindent\hangindent=2.4em\hangafter=1%
 \makebox[2.4em][l]{\sffamily\bfseries\color{Blue}#1}{\sffamily\bfseries #2}%
 \enspace{\color{Slate}#3}\par\vspace{1.1mm}}
% Contents at a glance.
\newenvironment{GlanceTable}{\par\small}{\par}
\newcommand{\GlancePart}[2]{\par\vspace{2.4mm}\noindent{\sffamily\bfseries\color{Blue}%
 Part \PartRoman{#1}\quad #2}\hfill{\sffamily\color{Blue}\pageref{book:part#1}}\par\vspace{0.6mm}}
\newcommand{\GlanceChapter}[3]{\noindent\hspace*{1.2em}\makebox[2.2em][l]{#1}#2\,\dotfill\,\pageref{#3}\par}
% Back matter chapter styles.
\newcommand{\EvolutorAppendixStyle}{%
 \titleformat{\chapter}[display]{\raggedright\hyphenpenalty=10000\color{Blue}\sffamily\bfseries}%
 {\large APPENDIX \thechapter}{8pt}{\huge}%
 \fancyhead[R]{\small\sffamily Appendix \thechapter}%
 \renewcommand{\thefigure}{\thechapter.\arabic{figure}}}
\newcommand{\UnnumberedChapter}[2]{\chapter*{#1}\addcontentsline{toc}{chapter}{#1}%
 \markboth{#1}{#1}\fancyhead[R]{\small\sffamily #2}}
\newcommand{\BackPlate}[3]{\par\medskip\noindent\begin{minipage}{\linewidth}\centering
\sbox{\PlateBox}{\input{convergence/figures/#1.tex}}\PlaceScaledPlate
\captionof{figure}{#2}\OriginalBookLabel{#3}\end{minipage}\par\medskip}
